% ===========================================================================
\section{子环、同态与理想}\label{sec:subring-hom-ideal}

\subsection{子环与环同态}

\begin{definition}[子环]
$R'\subseteq R$ 称为 $R$ 的\keyword{子环}，若它在同样的加法与乘法下
自身构成环。验证时只需检查 $R'$ 对\keyword{减法与乘法封闭}。
\end{definition}

\begin{definition}[环同态]
映射 $f\colon R\to R'$ 称为\keyword{（环）同态}，若对任意 $a,b\in R$：
\[
  f(a+b) = f(a) + f(b), \qquad f(ab) = f(a)f(b).
\]
由于 $f$ 保加法，它自动是加法群同态，故 $f(0) = 0$，$f(-a) = -f(a)$。
\end{definition}

\begin{proposition}[像是子环]
同态 $f\colon R\to R'$ 的像 $\Img f = \{\,f(a) : a\in R\,\}$ 是 $R'$
的子环。
\end{proposition}

\begin{proof}
由群论已知 $\Img f$ 对减法封闭；又 $f(a)f(b) = f(ab)\in\Img f$，
对乘法封闭。
\end{proof}

\subsection{核与理想}

\begin{definition}[核]
环同态 $f\colon R\to R'$ 的\keyword{核}为
\[
  \ker f = \{\,b\in R : f(b) = 0\,\}.
\]
\end{definition}

核是子环，且还有更强的性质：对任意 $a\in R$、$b\in\ker f$，
\[
  f(ab) = f(a)f(b) = f(a)\cdot 0 = 0 \;\Longrightarrow\; ab\in\ker f,
\]
同理 $ba\in\ker f$。即\keyword{环中任意元素乘核中元素（无论左右）仍
落在核中}。

\begin{definition}[理想（Ideal）]
环 $R$ 的非空子集 $I$ 称为 $R$ 的\keyword{理想}，记 $I\ideal R$，若
\begin{enumerate}[leftmargin=1.6em,itemsep=2pt]
  \item $(I,+)$ 是 $(R,+)$ 的子群（对减法封闭）；
  \item \keyword{吸收性}：对任意 $a\in R$、$b\in I$，有 $ab\in I$ 且
        $ba\in I$。
\end{enumerate}
\end{definition}

\begin{strand}
理想的定义正是\keyword{照着核的性质}写出来的，因此同态的核自然是
理想——理想在环论中的地位对应正规子群在群论中的地位。
\end{strand}

\subsection{商环（Quotient Ring）}

设 $I$ 是环 $R$ 的理想。$(I,+)$ 是阿贝尔群 $(R,+)$ 的子群（自动正规），
可考虑加法陪集
\[
  a + I = \{\,a + i : i\in I\,\},
\]
其全体记为 $R/I$。在 $R/I$ 上定义运算：
\[
  (a+I) + (b+I) = (a+b) + I, \qquad (a+I)(b+I) = ab + I.
\]

\begin{theorem}[商环]
$(R/I, +, \cdot)$ 构成环，称为 $R$ 模 $I$ 的\keyword{商环}。
\end{theorem}

\begin{proof}
加法部分的全部性质由商群理论保证；乘法只需验证\keyword{良好定义性}。
设 $a + I = a' + I$，$b + I = b' + I$，则 $a' = a + i$，$b' = b + j$
（$i,j\in I$），于是
\[
  a'b' = (a+i)(b+j) = ab + aj + ib + ij.
\]
由吸收性，$aj, ib, ij\in I$，故 $aj + ib + ij\in I$，从而
\[
  a'b' + I = ab + I,
\]
即乘积与代表元选取无关。
\end{proof}
